\documentclass[aspectratio=169, 10pt, xcolor=dvipsnames]{beamer}
\usepackage[UTF8]{ctex}
\usepackage{graphicx}
\usepackage{booktabs}
\usepackage{array}
\usepackage{colortbl}
\usepackage{hyperref}
\usepackage{fontawesome5}

% ── Theme ────────────────────────────────────────────────────────────────
\usetheme{Madrid}
\usecolortheme{whale}
\setbeamertemplate{navigation symbols}{}
\setbeamertemplate{footline}[frame number]
\setbeamerfont{frametitle}{size=\large, series=\bfseries}

% ── Colors ───────────────────────────────────────────────────────────────
\definecolor{healthgreen}{RGB}{34,139,34}
\definecolor{overfitred}{RGB}{200,30,30}
\definecolor{trainblue}{RGB}{30,90,200}
\definecolor{lightgray}{RGB}{240,240,240}
\definecolor{bestrow}{RGB}{220,245,220}

% ── Macros ───────────────────────────────────────────────────────────────
\newcommand{\imgwidth}{0.88\textwidth}
\graphicspath{{sources/}}

% ── Title ────────────────────────────────────────────────────────────────
\title[Z1 训练分析]{Z1 12DOF 学习曲线分析}
\subtitle{12 轮训练的 Reward 对比、Reward 分解、终止原因与训练效率}
\author{MagicLab RL Lab}
\date{2026-05-05}
\institute{RTX 6000D $\times$ 4 \quad | \quad Isaac Lab 0.47.2}

% ==========================================================================
\begin{document}

% ── Title page ───────────────────────────────────────────────────────────
\begin{frame}[plain]
  \titlepage
\end{frame}

% ── Outline ──────────────────────────────────────────────────────────────
\begin{frame}{目录}
  \tableofcontents
\end{frame}

% ==========================================================================
\section{训练概览}
% ==========================================================================

\begin{frame}{训练概览}
  \begin{block}{实验配置}
    \small
    硬件: 4$\times$ RTX 6000D (85.7 GB) \quad|\quad
    框架: Isaac Lab 0.47.2 + Isaac Sim 4.5.0 \\[2pt]
    任务: Z1 12DOF Velocity Locomotion \quad|\quad
    数据源: TensorBoard + \texttt{best\_models.json}
  \end{block}

  \vspace{0.3cm}

  \begin{itemize}
    \item 共 \textbf{12 个有效 run}（排除多 GPU 测试和重复启动）
    \item 按最佳模型 Reward 降序排列
    \item \textcolor{healthgreen}{\textbf{HEALTHY}}: 训练正常 \quad
          \textcolor{overfitred}{\textbf{OVERFITTING}}: reward 从 peak 大幅下降 \quad
          \textcolor{trainblue}{\textbf{TRAINING}}: 训练中
  \end{itemize}
\end{frame}

% ── Overview table (split into 2 slides) ─────────────────────────────────

\begin{frame}{训练概览 — 高 Reward 组}
  \centering\scriptsize
  \begin{tabular}{>{\bfseries}l l r r r l}
    \toprule
    Run & 别名 & 迭代数 & Peak & Best (iter) & 状态 \\
    \midrule
    \rowcolor{bestrow}
    v1\_flat & z1\_v1\_flat & 36,429 & 48.35 & 47.33 (3,861) & \textcolor{overfitred}{OVERFITTING} \\
    \rowcolor{bestrow}
    v2\_flat\_retry & z1\_v2\_retry & 31,173 & 48.35 & 47.33 (3,861) & \textcolor{overfitred}{OVERFITTING} \\
    \rowcolor{bestrow}
    v4\_gentle & z1\_v4\_gentle & 49,999 & 48.10 & \textbf{47.06} (47,862) & \textcolor{healthgreen}{\textbf{HEALTHY}} \\
    v5\_rough & z1\_v5\_rough & 41,590 & 38.94 & 38.04 (32,790) & \textcolor{overfitred}{OVERFITTING} \\
    v6\_l1\_4gpu & v6\_l1\_4gpu & 7,735 & 33.28 & 31.20 (5,032) & \textcolor{overfitred}{OVERFITTING} \\
    v3\_highspeed & z1\_v3 & 3,812 & 30.87 & 30.11 (2,997) & \textcolor{overfitred}{OVERFITTING} \\
    \bottomrule
  \end{tabular}

  \vspace{0.4cm}
  \small\textcolor{healthgreen}{\faCheckCircle} \textbf{v4\_gentle} 是唯一健康完成、peak reward 达到 48+ 的 run
\end{frame}

\begin{frame}{训练概览 — 低 Reward 及进行中}
  \centering\scriptsize
  \begin{tabular}{>{\bfseries}l l r r r l}
    \toprule
    Run & 别名 & 迭代数 & Peak & Best (iter) & 状态 \\
    \midrule
    v2\_stable & z1\_v2\_stable & 12,834 & 29.99 & 28.93 (1,555) & \textcolor{overfitred}{OVERFITTING} \\
    v6\_l1 & v6\_l1 & 1,799 & 7.37 & 5.86 (1,778) & \textcolor{overfitred}{OVERFITTING} \\
    v4\_terrain & z1\_v4\_terrain & 3,143 & 2.86 & 1.85 (1,933) & \textcolor{overfitred}{OVERFITTING} \\
    v6 (4gpu) & z1\_mgpu\_v6 & 50,018 & -0.46 & -0.46 (49,999) & \textcolor{overfitred}{OVERFITTING} \\
    v6\_4gpu\_s4 & z1\_mgpu\_4gpu & 50,070 & -1.68 & -1.68 (49,999) & \textcolor{overfitred}{OVERFITTING} \\
    \midrule
    \textcolor{trainblue}{\textbf{s4\_full}} & \textcolor{trainblue}{s4\_full} & \textcolor{trainblue}{训练中} & \textcolor{trainblue}{TBD} & \textcolor{trainblue}{TBD} & \textcolor{trainblue}{\textbf{TRAINING}} \\
    \bottomrule
  \end{tabular}

  \vspace{0.4cm}
  \small 多 GPU runs (v6 系列) 始终未能学到有效策略，reward 为负值
\end{frame}

% ==========================================================================
\section{Reward 对比分析}
% ==========================================================================

\begin{frame}{Plot 1: Reward 对比}
  \centering
  \includegraphics[width=\imgwidth]{1_reward_comparison.png}
\end{frame}

\begin{frame}{Reward 对比分析}
  \begin{columns}[T]
    \begin{column}{0.48\textwidth}
      \textbf{表现最好}
      \begin{itemize}
        \item \textbf{v1\_flat / v2\_flat\_retry} \\
              Peak $\sim$48.3，但 $\sim$3k iter 后崩溃
        \item \textbf{v4\_gentle} \\
              Peak $\sim$48.1，50k iter 仍保持 $\sim$45 \\
              \textcolor{healthgreen}{唯一健康完成的 run}
      \end{itemize}
    \end{column}
    \begin{column}{0.48\textwidth}
      \textbf{收敛与过拟合}
      \begin{itemize}
        \item v1/v2/v3: 快速收敛但 policy collapse
        \item v4\_gentle: 慢而稳（$\sim$48k iter 达 peak）
        \item v5\_rough: 32.8k 达 peak 后崩溃
        \item 多 GPU runs: 始终负 reward
      \end{itemize}
    \end{column}
  \end{columns}

  \vspace{0.3cm}
  \begin{alertblock}{关键观察}
    早期 flat terrain runs 达到 peak 后 reward 骤降至负值 — 典型的 \textbf{policy collapse} 模式。温和地形课程是唯一能稳定收敛的策略。
  \end{alertblock}
\end{frame}

% ==========================================================================
\section{Reward 分解}
% ==========================================================================

\begin{frame}{Plot 2: Reward 分解 (s4\_full)}
  \centering
  \includegraphics[width=\imgwidth]{2_reward_decomposition_s4_full.png}
\end{frame}

\begin{frame}{Reward 分解分析}
  \begin{columns}[T]
    \begin{column}{0.48\textwidth}
      \textbf{主要奖励组件}
      \vspace{0.2cm}

      \scriptsize
      \begin{tabular}{lll}
        \toprule
        组件 & 贡献 & 说明 \\
        \midrule
        +tracking\_lin\_vel & \textcolor{healthgreen}{\textbf{正向主导}} & 线速度跟踪 \\
        +tracking\_ang\_vel & \textcolor{healthgreen}{正向} & 角速度跟踪 \\
        +alive & \textcolor{healthgreen}{小正向} & 存活奖励 \\
        \midrule
        -action\_rate & \textcolor{overfitred}{负向} & 动作平滑惩罚 \\
        -torques & \textcolor{overfitred}{小负向} & 力矩惩罚 \\
        -dof\_vel & \textcolor{overfitred}{小负向} & 关节速度惩罚 \\
        \bottomrule
      \end{tabular}
    \end{column}
    \begin{column}{0.48\textwidth}
      \textbf{Curriculum 进度}
      \begin{itemize}
        \item Terrain level 从 0 逐步上升
        \item Velocity cmd level 同步增加
        \item 渐进式课程使 policy 逐步适应更复杂环境
      \end{itemize}
    \end{column}
  \end{columns}
\end{frame}

% ==========================================================================
\section{终止原因分析}
% ==========================================================================

\begin{frame}{Plot 3: 终止原因 (s4\_full)}
  \centering
  \includegraphics[width=\imgwidth]{3_termination_s4_full.png}
\end{frame}

\begin{frame}{终止原因分析}
  \begin{columns}[T]
    \begin{column}{0.48\textwidth}
      \textbf{终止原因趋势}
      \begin{itemize}
        \item \textcolor{overfitred}{\textbf{bad\_orientation}}: 初期 $\sim$100\% → 逐步下降
        \item \textcolor{healthgreen}{\textbf{time\_out}}: 逐步上升，episode 成功完成
        \item \textcolor{orange}{\textbf{base\_height}}: 偶发
      \end{itemize}
    \end{column}
    \begin{column}{0.48\textwidth}
      \textbf{Episode 长度}
      \begin{itemize}
        \item 初期很短（robot 立即摔倒）
        \item 逐渐接近上限 1000 steps (20s)
        \item 理想状态: time\_out $>$80\%
      \end{itemize}
    \end{column}
  \end{columns}

  \vspace{0.3cm}
  \begin{exampleblock}{训练进步指标}
    \texttt{time\_out} 比例上升 + \texttt{bad\_orientation} 比例下降 = robot 学会了保持平衡。
    训练后期目标: time\_out $\geq$ 80\%, bad\_orientation $\leq$ 10\%.
  \end{exampleblock}
\end{frame}

% ==========================================================================
\section{训练效率}
% ==========================================================================

\begin{frame}{Plot 4: 训练效率 (s4\_full)}
  \centering
  \includegraphics[width=\imgwidth]{4_efficiency_s4_full.png}
\end{frame}

\begin{frame}{训练效率分析}
  \begin{columns}[T]
    \begin{column}{0.48\textwidth}
      \textbf{Throughput}
      \begin{itemize}
        \item 单 GPU: $\sim$176k steps/s
        \item 4 GPU: $\sim$524k steps/s (3$\times$)
        \item 8$\times$ envs $\to$ 3$\times$ throughput
        \item PPO 同步是瓶颈
      \end{itemize}

      \vspace{0.2cm}
      \textbf{Entropy}
      \begin{itemize}
        \item 初始 $\sim$4.0+，逐步下降
        \item 过早 $\to$ 0 = 探索不足
      \end{itemize}
    \end{column}
    \begin{column}{0.48\textwidth}
      \textbf{Collection vs Learning}
      \begin{itemize}
        \item Collection: 仿真耗时（主导）
        \item Learning: PPO 更新耗时
        \item envs $\uparrow$ $\to$ learning time $\uparrow\uparrow$
      \end{itemize}

      \vspace{0.2cm}
      \textbf{Learning Rate}
      \begin{itemize}
        \item Adaptive schedule (rsl-rl)
        \item 随 KL divergence 自适应调整
        \item 训练后期逐步衰减
      \end{itemize}
    \end{column}
  \end{columns}
\end{frame}

% ==========================================================================
\section{关键发现与建议}
% ==========================================================================

\begin{frame}{关键发现}
  \begin{enumerate}
    \item \textbf{Flat terrain 训练不稳定} \\
          v1\_flat / v2\_flat\_retry 达到 $\sim$48 reward 后崩溃，缺乏多样性导致过拟合
    \item \textbf{温和地形课程最有效} \\
          v4\_gentle: 50k iter 达 peak 47.06 且保持稳定 — \textcolor{healthgreen}{\textbf{唯一健康完成的 run}}
    \item \textbf{Rough terrain 过早崩溃} \\
          v5\_rough 在 32.8k iter 达 peak 38.9 后崩溃，terrain 难度上升过快
    \item \textbf{L1 action rate 惩罚效果有限} \\
          v6\_l1 peak 仅 7.37
    \item \textbf{多 GPU 未改善训练质量} \\
          v6/v6\_4gpu reward 为负，配置问题
  \end{enumerate}
\end{frame}

\begin{frame}{最佳 Checkpoint 推荐}
  \centering
  \begin{tabular}{lll}
    \toprule
    \textbf{用途} & \textbf{推荐 Checkpoint} & \textbf{备注} \\
    \midrule
    \rowcolor{bestrow}
    部署 / 仿真验证 & v4\_gentle / model\_47862 (47.06) & 最稳定，50k iter \\
    Flat 基准对比 & v1\_flat / model\_3861 (47.33) & Reward 最高但需注意崩溃 \\
    Rough terrain & v5\_rough / model\_32790 (38.04) & 唯一 rough terrain 数据点 \\
    \bottomrule
  \end{tabular}
\end{frame}

\begin{frame}{下一步训练建议}
  \begin{columns}[T]
    \begin{column}{0.48\textwidth}
      \textbf{立即行动}
      \begin{enumerate}
        \item 继续 s4\_full 训练 \\
              (32768 envs / 4 GPU, $\sim$55k iter)
        \item 参考 v4\_gentle 的 curriculum 节奏 \\
              放缓 terrain 难度提升
        \item 监控 action\_rate $>$ -1.0
      \end{enumerate}
    \end{column}
    \begin{column}{0.48\textwidth}
      \textbf{策略调整}
      \begin{enumerate}
        \item Early stopping: 连续 5k iter \\
              reward 下降 $>$20\% 则停止
        \item 32768 envs 时 max\_iterations \\
              按比例缩小到 $\sim$14k
        \item 排查多 GPU 配置问题
      \end{enumerate}
    \end{column}
  \end{columns}

  \vspace{0.5cm}
  \centering
  \begin{block}{}
    \centering
    \small 使用 \texttt{/plot-train-Z1} 命令随时更新学习曲线图和分析文档
  \end{block}
\end{frame}

% ── Thank you ────────────────────────────────────────────────────────────
\begin{frame}[plain]
  \centering
  \vfill
  {\Huge\bfseries 谢谢} \\[1cm]
  {\large 问题与讨论} \\[0.5cm]
  {\small\textcolor{gray}{生成工具: plot\_learning\_curves.py + LaTeX Beamer}}
  \vfill
\end{frame}

\end{document}
